\chapter{共形映射与黎曼映射定理}
\label{ch:b3:conformal}

两个区域之间的全纯双射，把复分析的全部内容
从一个区域搬到另一个区域：这样的映射 ---
\emph{共形}的，因其保持角度 --- 是
全纯世界的同构。本章对它们作分类，
在可以分类之处（单位圆盘与复平面：施瓦茨引理是关键，
这是一个威力惊人的不等式），
并建立全纯函数族的紧性理论
（蒙特尔），并证明本学科最深刻的存在性定理：
\emph{每一}个
$\C$ 的单连通真子区域，无论其边界多么参差，都与
单位圆盘共形等价。本章以调和函数与
泊松核作结，求解圆盘上的狄利克雷问题 ---
共形几何的分析回报。全文中，
$\mathbb D = D(0,1)$ 与 $\mathbb H = \{\operatorname{Im}z >
0\}$。

\section{共形映射；莫比乌斯变换}

\begin{definition}\label{def:b3:conformal:conformal}
称\emph{共形映射}\index{共形映射}（或
双全纯映射）是开集之间的全纯双射；
其逆映射自动全纯
（见\cref{cor:b3:residues:openmapping}）。两个区域若
存在此类映射，则称二者为\emph{共形等价}；
$\operatorname{Aut}(\Omega)$ 表示由共形
自映射构成的群。当 $f' \neq 0$ 时 ---
对单射 $f$ 处处成立（见\cref{cor:b3:residues:openmapping}的证明）--- 其
微分即乘以 $f'(z) \neq 0$：一个相似变换，
因此共形映射保持曲线之间的夹角，
包括定向。
\end{definition}

\begin{example}[莫比乌斯变换]
\label{ex:b3:conformal:mobius}
设 $\bigl(\begin{smallmatrix}a & b\\ c &
d\end{smallmatrix}\bigr) \in GL_2(\C)$，则\emph{莫比乌斯
变换}\index{莫比乌斯变换} $z \mapsto
\frac{az + b}{cz + d}$ 是从
$\C\setminus\{-d/c\}$ 到 $\C\setminus\{a/c\}$ 上的共形映射（逆
映射同属此类，由逆矩阵给出；复合
对应于矩阵乘积）。\emph{凯莱
映射}\index{凯莱映射}
\[
\varphi(z) = \frac{z - \iu}{z + \iu}
\]
将 $\mathbb H$ 共形地映为 $\mathbb D$，事实上 $\abs{z
- \iu} < \abs{z + \iu}$ 当且仅当 $z$ 比较接近 $\iu$
而不是 $-\iu$，即 $\operatorname{Im}z > 0$；其逆映射为
$w \mapsto \iu\frac{1 + w}{1 - w}$。莫比乌斯变换把
圆与直线构成的族映为自身（见\cref{exo:b3:conformal:1}）。
\end{example}

\begin{example}[儒可夫斯基映射]\label{ex:b3:conformal:joukowski}
除莫比乌斯变换外，经典应用数学中最有用的
共形映射是
\[
J(z) = \frac12\Bigl(z + \frac1z\Bigr) .
\]
在单位圆盘的外部 $\Omega = \{\abs z > 1\}$ 上，
$J$ 是单射：由 $J(z) = J(w)$ 得 $(z - w)(1 -
\frac1{zw}) = 0$，且 $\abs{zw} > 1$。其导数 $J'(z) =
\frac12(1 - z^{-2})$ 仅在 $z = \pm1$ 处为零，这两点
位于边界上：$J$ 是 $\Omega$ 到其像上的共形等价，
其像即 $\C\setminus\intcc{-1}1$ --- 单位
圆周本身被二对一地折叠到线段上
（$J(\eu^{\iu\theta}) = \cos\theta$）。因此，一条
\emph{线段}的外部，一个没有光滑边界的缝隙平面，
在共形意义下即圆盘的外部：角点并不妨碍
共形等价，只妨碍边界的光滑性。经过
$\pm1$ 而偏离中心的圆之像是翼型状
曲线；将 $J$ 与莫比乌斯变换复合，便把
绕圆柱的流动（可以手算）--- 变为
绕机翼的流动：二十世纪上半叶，
这一例子\emph{正是}空气动力学。它也是通向
切比雪夫理论的门径：$J$ 将 $z \mapsto z^n$ 与切比雪夫
多项式 $T_n$ 联系起来（见\cref{pb:b3:hilbert:1}，V 部分），因为
$J(z^n) = \cos(n\theta)$，当 $z = \eu^{\iu\theta}$ 时。
\end{example}

\section{施瓦茨引理与自同构群}

\begin{theorem}[施瓦茨引理]\label{thm:b3:conformal:schwarz}
设 $f \colon \mathbb D \to \mathbb D$ 为全纯函数，且
$f(0) = 0$。则\index{施瓦茨引理}
\[
\abs{f(z)} \leq \abs z \ \ (z \in \mathbb D)
\qquad\text{且}\qquad \abs{f'(0)} \leq 1 ,
\]
若 $\abs{f(z_0)} = \abs{z_0}$ 对某个 $z_0 \neq 0$ 成立，或
$\abs{f'(0)} = 1$，则 $f(z) = \eu^{\iu\theta}z$ 是一个
旋转。
\end{theorem}

\begin{proof}
$g(z) = f(z)/z$ 可全纯地延拓到 $\mathbb D$ 上（
$0$ 处的奇点可去：$g$ 在 $0$ 附近有界，
见\cref{thm:b3:residues:riemanncw}；其在 $0$ 处的值为
$f'(0)$）。在 $\abs z = r < 1$ 上，$\abs g \leq \frac1r$，由
最大模原理（见\cref{thm:b3:holomorphic:maximum}）得
$\abs g \leq \frac1r$ 在 $\bar D(0,r)$ 上；令 $r \to 1$，$\abs
g \leq 1$ 在 $\mathbb D$ 上，这即两个不等式。等号
在内点处成立时，$\abs g$ 在该内点取得最大值，
从而 $g$ 为模等于 $1$ 的常数。
\end{proof}

\begin{theorem}[圆盘的自同构]
\label{thm:b3:conformal:autdisc}
对 $a \in \mathbb D$，\emph{布拉施克
因子}\index{布拉施克因子}
\[
\varphi_a(z) = \frac{z - a}{1 - \bar a z}
\]
是 $\mathbb D$ 的一个自同构，交换 $a$ 与 $0$，并有
$\varphi_a^{-1} = \varphi_{-a}$。$\mathbb D$ 的每个自同构
都可唯一地写成 $\eu^{\iu\theta}\varphi_a$，其中 $\theta
\in \R/2\pi\Z$，$a \in \mathbb D$。
\end{theorem}

\begin{proof}
在 $\abs z = 1$ 上：$\abs{1 - \bar az} = \abs{\bar z}\abs{1 -
\bar az} = \abs{\bar z - \bar a\abs z^2} = \abs{\bar z - \bar
a} = \abs{z - a}$，故 $\abs{\varphi_a} = 1$ 于此；由
最大模原理得 $\varphi_a(\mathbb D) \subseteq \bar{\mathbb
D}$，由开性，像落在 $\mathbb D$ 内。代数
恒等式 $\varphi_{-a}\circ\varphi_a = \mathrm{id}$（直接
计算）可知其为双射。设 $f \in
\operatorname{Aut}(\mathbb D)$，$a = f^{-1}(0)$，记 $g =
f\circ\varphi_{-a}$ 是固定 $0$ 的自同构。施瓦茨
引理用于 $g$ 与 $g^{-1}$：$\abs{g(z)} \leq \abs z$ 及
$\abs{g^{-1}(w)} \leq \abs w$，故 $\abs{g(z)} = \abs z$：
旋转，$g = \eu^{\iu\theta}\,\mathrm{id}$，即 $f =
\eu^{\iu\theta}\varphi_a$。唯一性：$a = f^{-1}(0)$ 及
$\theta$ 可由 $f'$ 型取值确定（或由 $f(0) =
-\eu^{\iu\theta}a$ 及另一取值）。
\end{proof}

\begin{theorem}[平面的自同构]
\label{thm:b3:conformal:autplane}
$\operatorname{Aut}(\C) = \{z \mapsto az + b : a \in \C^*,\ b
\in \C\}$。因此 $\C$ 与 $\mathbb D$ 不
共形等价。
\end{theorem}

\begin{proof}
设 $f \in \operatorname{Aut}(\C)$，考虑 $g(z) =
f(1/z)$ 在 $\C^*$ 上，是一个具有孤立
奇点 $0$ 的全纯函数。若它为本质奇点，
则由卡索拉蒂–魏尔斯特拉斯定理（\cref{thm:b3:residues:riemanncw}）将使
$g\bigl(D(0,\varepsilon)\setminus\{0\}\bigr)$ 稠密，
而 $f(D(0, 1))$ 为开集，且与之不交（$f$
为单射：两集合是不交集合的像） ---
这对稠密集与非空开集不可能。故 $0$ 是
$g$ 的极点或可去奇点，即 $\abs{f(z)}$ 至多
多项式增长：$f$ 为多项式
（见\cref{exo:b3:holomorphic:4}(a)）。单射性要求次数
$1$：次数更高的多项式要么有某个 $f - c$ 的重根
（即某处 $f'$ 为零），要么有若干个相异
原像（达朗贝尔–高斯，见\cref{pb:b3:holomorphic:1}）；
两种情形下单射性均不成立。最后，若存在共形映射 $\C \to
\mathbb D$，则其为有界整函数：
由刘维尔定理可知它为常值 --- 无共形等价。
\end{proof}

\section{蒙特尔定理}

\begin{theorem}[蒙特尔]\label{thm:b3:conformal:montel}
设 $\mathcal F \subseteq \mathcal H(\Omega)$ 为\emph{局部
有界}：每点都有邻域，在其上 $\sup_{f\in
\mathcal F}\sup\abs f < \infty$。则
$\mathcal F$ 的每个序列都有子列，在所有
$\Omega$ 的紧子集上一致收敛（收敛于全纯的
极限）。\index{蒙特尔定理}
\end{theorem}

\begin{proof}
局部等度连续：若 $\abs f \leq M$ 于 $D(a, 2r)
\subseteq \Omega$，对所有 $f \in \mathcal F$，由柯西
公式给出，对 $z, z' \in D(a, r)$，
\[
\abs{f(z) - f(z')} = \frac{\abs{z - z'}}{2\pi}
\Bigl|\int_{C_{2r}}\frac{f(w)\,\dd w}{(w-z)(w-z')}\Bigr|
\leq \frac{\abs{z - z'}\,2\pi\cdot2r\,M}{2\pi\,r^2}
= \frac{2M}{r}\,\abs{z - z'} :
\]
即一致利普希茨界。将 $\Omega$ 以紧集 $K_m$ 穷竭；
每个 $K_m$ 被有限个此类圆盘覆盖，故
$\mathcal F$ 在
$K_m$ 上一致有界且等度连续：由阿尔泽拉–阿斯科利（\cref{thm:b3:complete:ascoli}）
可取出在 $K_m$ 上一致收敛的子列；
对 $m$ 作对角线法。由
\cref{thm:b3:holomorphic:weierstrassconv}可知极限为全纯函数。
\end{proof}

\section{黎曼映射定理}

\begin{definition}\label{def:b3:conformal:simplyconnected}
称开连通集 $\Omega \subseteq \C$ 为\emph{单
连通}\index{单连通}（在同调意义下，
足以满足本课程全部用途）若
$\operatorname{Ind}_\gamma(w) = 0$，对每条闭链 $\gamma$
（在 $\Omega$ 中）及每个 $w \notin \Omega$ 成立 --- “$\Omega$ 中
没有闭链环绕孔洞”。由全局柯西定理
（\cref{thm:b3:residues:globalcauchy}）及
\cref{prop:b3:holomorphic:primitive}，在此类 $\Omega$ 上
\emph{每个全纯函数都有原函数}；因此每个
无零点的 $f \in \mathcal H(\Omega)$ 有全纯
对数（$\exp\circ$（$f'/f$ 的原函数），经
常数修正，因 $(f\eu^{-L})' = 0$），并有全纯 $n$ 次方根
$\eu^{L/n}$。\index{全纯对数}
\end{definition}

\begin{theorem}[黎曼映射定理]
\label{thm:b3:conformal:rmt}
每个单连通开集 $\Omega \subsetneq \C$，$\Omega
\neq \varnothing$，与 $\mathbb D$ 共形等价；
给定 $z_0 \in \Omega$，存在唯一的共形映射 $f \colon
\Omega \to \mathbb D$，使 $f(z_0) = 0$ 且 $f'(z_0) >
0$。\index{黎曼映射定理}
\end{theorem}

\begin{proof}
\emph{第 0 步：族非空。}取 $b \notin \Omega$：
$z - b$ 在 $\Omega$ 上无零点，故其有全纯
平方根 $h$（$h^2 = z - b$）。$h$ 是单射（$h(z) =
h(z')$ 平方得 $z = z'$），且若 $w \in h(\Omega)$，则 $-w
\notin h(\Omega)$（$h(z) = -h(z')$ 平方同样得 $z = z'$，
从而 $w = -w = 0$，与 $h$ 无零点矛盾）。由于
$h(\Omega)$ 为开集，故含圆盘 $D(h(z_0), \rho)$；
则 $D(-h(z_0), \rho) \cap h(\Omega) = \varnothing$，即
$\abs{h(z) + h(z_0)} \geq \rho$ 对每个 $z \in \Omega$ 成立。
由此
\[
g(z) = \frac{\rho}{2\,\bigl(h(z) + h(z_0)\bigr)}
\]
全纯，且为单射（一个莫比乌斯映射与
单射 $h$ 复合），且 $\abs g \leq \frac12 < 1$。与布拉施克因子
复合（见\cref{thm:b3:conformal:autdisc}），使 $g(z_0)$ 移到 $0$，
该族
\[
\mathcal F = \{f \colon \Omega \to \mathbb D \text{
全纯、单射， } f(z_0) = 0\}
\]
非空。

\emph{第 1 步：极值元。}设 $s =
\sup_{\mathcal F}\abs{f'(z_0)} \in \intoc0{+\infty}$（$> 0$：
成员都是单射，故 $f'(z_0) \neq 0$）。取 $f_n \in
\mathcal F$，使 $\abs{f_n'(z_0)} \to s$：该族
以 $1$ 为界，故由蒙泰尔定理（见\cref{thm:b3:conformal:montel}）
可取出子列，使 $f_n \to f$ 于紧集上一致收敛；$f$
全纯，$f(z_0) = 0$，$\abs{f'(z_0)} = s$
（导数部分由\cref{thm:b3:holomorphic:weierstrassconv}
可得），特别地 $f$ 非常值；$f$
是单射（由赫尔维茨定理，见\cref{exo:b3:residues:8}(b)），且
$f(\Omega) \subseteq \bar{\mathbb D}$，故 $\subseteq
\mathbb D$（开映射）。于是 $f \in \mathcal F$ 取到
上确界：$s < \infty$。

\emph{第 2 步：极值映射是满射。}设 $a \in
\mathbb D\setminus f(\Omega)$。布拉施克变换
$\varphi_a\circ f$ 在单连通的
$\Omega$ 上无零点，故有全纯平方根 $F$（且
$F(\Omega) \subseteq \mathbb D$，因 $\abs F^2 =
\abs{\varphi_a\circ f} < 1$），单射（平方可
区分）。归一化：$G = \varphi_{F(z_0)}\circ F \in
\mathcal F$。还原得 $f = \varphi_{-a}\circ s_2 \circ
\varphi_{-F(z_0)}\circ G$，其中 $s_2(w) = w^2$；映射 $\Psi
= \varphi_{-a}\circ s_2\circ\varphi_{-F(z_0)} \colon \mathbb D
\to \mathbb D$ 全纯，$\Psi(0) = f(z_0) = 0$，且
\emph{不}是旋转（它不是单射，$s_2$ 亦不是）。
施瓦茨引理（严格情形）：$\abs{\Psi'(0)} < 1$，且由
链式法则 $f = \Psi\circ G$ 得 $\abs{f'(z_0)} =
\abs{\Psi'(0)}\,\abs{G'(z_0)} < \abs{G'(z_0)}$ ---
与极大性矛盾（注意 $G \in \mathcal F$）。从而 $f$
为满射，即共形等价。

\emph{第 3 步：归一化与唯一性。}将 $f$ 乘以
$\eu^{-\iu\arg f'(z_0)}$ 使 $f'(z_0) > 0$（这仍在
$\mathcal F$）。若 $f_1, f_2$ 均满足要求，则 $\psi =
f_2\circ f_1^{-1} \in \operatorname{Aut}(\mathbb D)$ 固定
$0$，且 $\psi'(0) = f_2'(z_0)/f_1'(z_0) > 0$；
由\cref{thm:b3:conformal:autdisc}，$\psi$ 为旋转
$\eu^{\iu\theta}$，且 $\eu^{\iu\theta} > 0$：$\psi =
\mathrm{id}$。
\end{proof}

\begin{remark}\label{rem:b3:conformal:rmtscope}
该定理是一个纯存在性结论，其适用范围之广令人惊叹
：正方形、半平面、挖去一条割缝的平面、
两个相切圆之间的区域、边界为分形的区域
--- 均与 $\mathbb D$ 共形等价。该定理不
给出任何公式（显式映射是例外：
\cref{exo:b3:conformal:5}），边界性态（一套更深的理论
--- 卡拉泰奥多里定理 --- 处理之），或
延拓到 $\C$ 或多连通区域时的唯一性：
圆环 $\{1 < \abs z < 2\}$ \emph{不}与
去心圆盘共形等价，而半径比不同的圆环彼此
不等价（这是一个确实更难的事实）。
\end{remark}

\section{调和函数与泊松核}

\begin{proposition}\label{prop:b3:conformal:harmonicholo}
设 $\Omega$ 单连通，$u \colon \Omega \to \R$
为调和函数\index{调和函数}（$\mathcal C^2$ 且
$\Delta u = u_{xx} + u_{yy} = 0$）。则 $u =
\operatorname{Re}F$，其中 $F$ 全纯，在相差
一个纯虚常数的意义下唯一。因此 $u$ 是 $\mathcal C^\infty$，
满足均值性质，且服从最大值
原理（除非为常数，否则不存在严格的内部极值）。
\end{proposition}

\begin{proof}
$g = u_x - \iu u_y$ 满足柯西–黎曼方程
（$P = u_x$，$Q = -u_y$：$P_x = u_{xx} = -u_{yy} = Q_y$，且
$P_y = u_{xy} = u_{yx} = -Q_x$），偏导数连续，故 $g$
全纯（见\cref{prop:b3:holomorphic:cauchyriemann}；
$\R$ 可微性由 $\mathcal C^1$ 推出）。设 $F_0$
为一个原函数（利用单连通性，
见\cref{def:b3:conformal:simplyconnected}）；则
$\operatorname{Re}F_0$ 的梯度为 $(u_x, u_y)$
（$F_0' = g$ 展开后，对
$F_0$ 用柯西–黎曼方程即得此式），故 $u - \operatorname{Re}F_0$ 为常数（$\Omega$
连通）：取 $F = F_0 + c$。性质
由 \cref{thm:b3:holomorphic:maximum} 及
\cref{exo:b3:holomorphic:10}（对最大值原理
应用于 $u$ 本身时，则同样取 $\eu^{F}$）。
\end{proof}

\begin{theorem}[泊松公式；圆盘上的
狄利克雷问题]\label{thm:b3:conformal:poisson}
对 $0 \leq r < 1$，定义\emph{泊松
核}\index{泊松核}
\[
P_r(\theta) = \sum_{n\in\Z}r^{\abs n}\eu^{\iu n\theta}
= \frac{1 - r^2}{1 - 2r\cos\theta + r^2} \;>\; 0 .
\]
设 $g \colon \partial\mathbb D \to \R$ 连续，且
对 $z = r\eu^{\iu\varphi} \in \mathbb D$，令
\[
u(z) = \frac1{2\pi}\int_0^{2\pi}
P_r(\varphi - t)\,g(\eu^{\iu t})\,\dd t .
\]
则 $u$ 在 $\mathbb D$ 上调和，且连续延拓
到 $\bar{\mathbb D}$ 上，边界值为 $g$：满足此条件的
唯一调和函数。\index{狄利克雷问题}
\end{theorem}

\begin{proof}
\emph{核恒等式}：对两个等比级数求和，
\[
\sum_{n\in\Z}r^{\abs n}\eu^{\iu n\theta}
= \operatorname{Re}\frac{1 + r\eu^{\iu\theta}}{1 -
r\eu^{\iu\theta}}
= \frac{1 - r^2}{\abs{1 - r\eu^{\iu\theta}}^2},
\]
即上式所示的商；正性显然，且
$\frac1{2\pi}\int_0^{2\pi}P_r = 1$（仅 $n = 0$ 一项保留）。

\emph{调和性}：取 $z = r\eu^{\iu\varphi}$，
\[
u(z) = \operatorname{Re}\biggl[\frac1{2\pi}\int_0^{2\pi}
\frac{\eu^{\iu t} + z}{\eu^{\iu t} - z}\,
g(\eu^{\iu t})\,\dd t\biggr],
\]
（方括号中的核的实部为 $P_r(\varphi - t)$：
计算），且方括号式关于 $z$ 在 $\mathbb D$ 上全纯
（见\cref{exo:b3:holomorphic:7}）：$u$ 是一个
全纯函数的实部，因此为调和函数。

\emph{边界值}：$P_r(\cdot)$ 是
当 $r \to 1^-$ 时的近似恒等式：总质量为 $1$，且对 $\delta \leq
\abs\theta \leq \pi$, $P_r(\theta) \leq \frac{1 - r^2}{1 -
2r\cos\delta + r^2} \to 0$，且一致成立。标准的分拆
（$g$ 在 $\eu^{\iu\varphi_0}$ 附近的连续性，
其余处的有界性）给出 $u(r\eu^{\iu\varphi}) \to
g(\eu^{\iu\varphi_0})$，当 $r\eu^{\iu\varphi} \to
\eu^{\iu\varphi_0}$ 时，且对边界点一致，
所得延拓是连续的。\emph{唯一性}：两解之差
在 $\mathbb D$ 上调和，在
闭包上连续，在边界上为零：由最大值原理
（作用于 $\pm$ 差）知差恒为零。
\end{proof}

\begin{omfigure}
\begin{tikzpicture}[scale=1.05]
  \begin{scope}
    \clip (-2.1,-2.1) rectangle (2.1,2.1);
    \foreach \c in {0.2, 0.5, 0.8, 1.1, 1.4}{
      \draw[omDef!60] plot[domain=-2.05:2.05, samples=80]
        (\x, {\c/(2*\x)});
      \draw[omDef!60] plot[domain=-2.05:2.05, samples=80]
        (\x, {-\c/(2*\x)});
      \draw[omThm!60, samples=80, domain=\c:2.9]
        plot ({sqrt((\x*\x-\c*\c)/1)/1}, {0}) ;
    }
    % hyperbolas x^2 - y^2 = c
    \foreach \c in {0.4, 1, 1.6}{
      \draw[omThm!70] plot[domain=-1.32:1.32, samples=60]
        ({cosh(\x)*sqrt(\c)}, {sinh(\x)*sqrt(\c)});
      \draw[omThm!70] plot[domain=-1.32:1.32, samples=60]
        ({-cosh(\x)*sqrt(\c)}, {sinh(\x)*sqrt(\c)});
    }
  \end{scope}
  \draw[->] (-2.2,0) -- (2.3,0);
  \draw[->] (0,-2.2) -- (0,2.3);
  \node[font=\small, omThm] at (1.75,-1.9)
    {$\operatorname{Re}z^2 = c$};
  \node[font=\small, omDef] at (-1.35,1.95)
    {$\operatorname{Im}z^2 = c$};
\end{tikzpicture}

{\small $\operatorname{Re}z^2 = x^2 - y^2$ 的等值线
（红色双曲线）与 $\operatorname{Im}z^2 = 2xy$（蓝色
双曲线）在 $0$ 以外处成直角相交：共形映射
（$z \mapsto z^2$，其中 $z \neq 0$）保持
坐标线的正交性。在 $z = 0$ 处，
导数为零，角度反而加倍。}
\end{omfigure}

\section{练习}

\begin{exercise}[$\star$]\label{exo:b3:conformal:1}
(a) 验证凯莱映射 $\varphi(z) = \frac{z -
\iu}{z + \iu}$ 是 $\mathbb H \to \mathbb D$ 的双射，且其逆映射即所给者，
并求 $\iu$、$0$、
$1$、$\infty$ 的像（取极限）。
(b) 证明 $z \mapsto 1/z$ 把圆和直线
映为圆和直线。\emph{（写出它们的公共方程
$\alpha\abs z^2 + \bar\beta z + \beta\bar z + \gamma = 0$,
$\alpha, \gamma \in \R$。）} 由此对所有 M\"obius
映射同样成立。
\end{exercise}

\begin{exercise}[$\star$]\label{exo:b3:conformal:2}
设 $f \colon \mathbb D \to \mathbb D$ 为全纯函数。
(a) 若 $f(0) = 0$，且 $f(a) = a$（某个 $a \neq 0$），证明 $f
= \mathrm{id}$.
(b) 若 $f$ 为自同构，在
$\mathbb D$ 中有两个不同的不动点，证明 $f = \mathrm{id}$
\emph{（用布拉施克因子作共轭，归结为 (a)）}。
\end{exercise}

\begin{exercise}[$\star\star$]\label{exo:b3:conformal:3}
（施瓦茨--皮克）设全纯函数 $f\colon \mathbb D \to
\mathbb D$，证明
\[
\frac{\abs{f'(z)}}{1 - \abs{f(z)}^2} \;\leq\;
\frac{1}{1 - \abs z^2}
\qquad (z \in \mathbb D),
\]
等号成立（在某一点处成立，从而处处成立）当且仅当 $f \in
\operatorname{Aut}(\mathbb D)$.
\emph{（对 $\varphi_{f(z)}\circ
f\circ\varphi_{-z}$ 应用施瓦茨引理。）} 解释：全纯自映射
收缩双曲度量。
\end{exercise}

\begin{exercise}[$\star\star$]\label{exo:b3:conformal:4}
求出显式的共形等价映射：
(a) 带形 $\{0 < \operatorname{Im}z < \pi\} \to \mathbb H$；
(b) 象限 $\{\operatorname{Re}z > 0,
\operatorname{Im}z > 0\} \to \mathbb H$;
(c) 半圆盘 $\mathbb D\cap\mathbb H \to$ 象限，
再 $\to \mathbb H$；
(d) 映射 $\mathbb D \to \mathbb D$ 将 $\frac12$ 映为 $0$，
且该点处的导数为正。
\end{exercise}

\begin{exercise}[$\star\star$]\label{exo:b3:conformal:5}
（a）证明不存在共形映射 $\C \to \mathbb D$ 或 $\C \to
\mathbb H$，也不存在 $\mathbb D \to \C$ 的共形映射。
（b）下列哪些与
$\mathbb D$ 共形等价？请借助\cref{thm:b3:conformal:rmt}或某种障碍
性论证说明理由：正方形；$\C\setminus\intoc{-\infty}0$；
$\mathbb D\setminus\{0\}$；$\{1 < \abs z < 2\}$。
\emph{（对于最后两个：去心圆盘的共形像
将借助
\cref{thm:b3:residues:riemanncw}（1）延拓过该去心点 --- 请展开论证。）}
\end{exercise}

\begin{exercise}[$\star\star$]\label{exo:b3:conformal:6}
设 $\mathcal F = \{f \in \mathcal H(\mathbb D) : f(0) = 1,\
\operatorname{Re}f > 0\}$.
（a）证明 $\mathcal F$ 局部有界。
\emph{（与凯莱型映射 $w \mapsto \frac{w -
1}{w + 1}$ 复合，该映射将右半平面映为 $\mathbb D$，再
应用施瓦茨引理。）}
（b）由此推出\emph{赫格洛茨界}：$\abs{f(z)} \leq
\frac{1 + \abs z}{1 - \abs z}$，对 $f \in \mathcal F$ 成立，
并讨论等号成立的可能性。
\end{exercise}

\begin{exercise}[$\star\star\star$]\label{exo:b3:conformal:7}
\cref{thm:b3:conformal:rmt}的证明在何处用到了每一个
假设？请追踪：（i）单连通性（用两次）；（ii）
$\Omega \neq \C$；（iii）连通性。再说明该定理
对 $\Omega = \C$ 以及圆环\emph{不成立}，
并指出证明中哪一步在每种情形下失效。
\end{exercise}

\begin{exercise}[$\star\star$]\label{exo:b3:conformal:8}
求 $\mathbb D$ 上在下列边界数据下的狄利克雷问题的解：
（a）$g(\eu^{\iu t}) = \cos t$；（b）$g(\eu^{\iu t}) =
\cos^2 t$；（c）$g = \mathbf 1_{\text{上半圆周}}$ ---
对（c）计算 $u(0)$，并借助均值
性质解释其意义。\emph{（将 $g$ 展开为傅里叶级数，并利用
$P_r$ 的级数：$u(r\eu^{\iu\varphi}) = \sum_n c_n(g)
r^{\abs n}\eu^{\iu n\varphi}$。）}
\end{exercise}

\begin{exercise}[$\star\star\star$]\label{exo:b3:conformal:9}
（哈纳克）设 $u \geq 0$ 为 $\mathbb D$ 上的调和函数，证明
对 $\abs z = r < 1$，有
\[
\frac{1 - r}{1 + r}\,u(0) \;\leq\; u(z) \;\leq\;
\frac{1 + r}{1 - r}\,u(0)
\]
\emph{（将泊松核夹在 $\frac{1-r}{1+r}$ 与
$\frac{1+r}{1-r}$ 之间；在稍小的圆盘上应用
表示公式并取极限）}。由此推出：
$\C$ 上有下界的调和函数为常数。
\end{exercise}

\begin{exercise}[$\star\star$]\label{exo:b3:conformal:10}
利用调和性的共形不变性（$u\circ f$ 在 $u$ 调和而 $f$ 全纯时
仍调和 --- 可借助
\cref{prop:b3:conformal:harmonicholo}在局部证明），求
上半平面上以
$\mathbf 1_{\intoo{-\infty}0}$ 为边界数据的狄利克雷问题的解，并证明
\[
u(x + \iu y) = \frac1\pi\,\arg(x + \iu y)
\qquad (\arg \in \intoo0\pi \text{ 在 } \mathbb H)
\]
它在 $\mathbb H$ 上调和（即某全纯对数的
虚部），并在每个 $x \neq
0$ 处满足所需边界极限，并经凯莱映射将其移到圆盘上，以重新推导
\cref{exo:b3:conformal:8}（c）。
\end{exercise}

\begin{exercise}[$\star\star$]\label{exo:b3:conformal:11}
（圆盘内的不动点与迭代）设 $f\colon\mathbb
D \to \mathbb D$ 为全纯函数。
（a）证明：若 $f$ 有 \emph{两个}不同的不动点，
则 $f = \mathrm{id}$ \emph{（借助自同构将其中一个移到 $0$，
再应用施瓦茨引理的等号情形）}。
（b）设 $f(0) = 0$ 且 $f$ 不是旋转。证明
迭代 $f^{\circ n} \to 0$ 在每个紧集
$\bar D(0, r)$（$r < 1$）上一致收敛。\emph{（由施瓦茨引理，$\abs{f(z)}
\leq c_r\abs z$ 在 $\bar D(0,r)$ 上成立，其中 $c_r < 1$ --- 请说明
这一严格常数的来由，借助最大模原理作用于
$f(z)/z$）}。
（c）以 $f(z) = \frac{z^2 + z}2$ 为例，求出不动点，
并考察 $z_0 =
\frac12$ 的轨道的收敛速度。
\end{exercise}

\begin{exercise}[$\star\star$]\label{exo:b3:conformal:12}
（具体的调和共轭）设 $u(x, y) = x^3 - 3xy^2
+ 2y$.
（a）验证 $u$ 在 $\R^2$ 上调和，并求出所有
调和共轭 $v$（即 $u + \iu v$ 全纯），通过
积分柯西–黎曼方程；将 $f(z) =
u + \iu v$ 识别为 $z$ 的多项式。
（b）证明在\emph{星形}开集上，每个
调和函数都存在调和共轭，且在相差一个
加性常数的意义下唯一\emph{（$1$-形式 $-u_y\,\dd x +
u_x\,\dd y$ 是闭的；可借助\cref{thm:b3:holomorphic:cauchy}的
原函数方法，或\cref{ch:b3:forms}中的庞加莱
引理）}。
（c）给出标准反例，
即 $\C^*$ 上的 $u = \ln\abs z$ 没有整体共轭 --- 请
联系角形式与环绕数
（见\cref{ch:b3:forms}）。
\end{exercise}

\section{问题：面积定理与科贝
四分之一定理}

\begin{problem}[{周末问题 --- 单叶映射必须覆盖
多大区域？}]\label{pb:b3:conformal:1}
全纯单射称为\emph{单叶函数}。
圆盘上的规范化单叶函数，
\[
\mathcal S = \bigl\{f \in \mathcal H(\mathbb D) \text{
单射},\ f(z) = z + a_2z^2 + a_3z^3 + \cdots\bigr\},
\]
受到严格约束：证明比伯巴赫不等式
$\abs{a_2} \leq 2$，并由此推出\emph{柯贝四分之一
定理}：任一 $f \in \mathcal S$ 的像都包含
圆盘 $D(0, \frac14)$ --- 共形几何中最优的普适常数
。\index{柯贝四分之一定理}\index{面积
定理}

\medskip
\noindent\textbf{第 I 部分 --- 面积定理。}设 $g(w) = w
+ b_0 + \frac{b_1}w + \frac{b_2}{w^2} + \cdots$ 全纯且
\emph{单射}，定义域为 $\{\abs w > 1\}$。
\begin{enumerate}
  \item 对 $\rho > 1$，记 $A_\rho$ 为紧集
        $K_\rho = \C\setminus
        g(\{\abs w > \rho\})$ 的面积（即勒贝格测度），它是由
        光滑若尔当曲线 $g(C_\rho)$ 所围的区域。利用本科二年级分册中的
        格林–黎曼定理的面积公式
        --- 所围面积为 $\frac1{2\iu}
        \oint\bar\zeta\,\dd\zeta$ 沿正向
        边界。证明：
        \[
        A_\rho = \frac{1}{2\iu}\int_{C_\rho}
        \overline{g(w)}\,g'(w)\,\dd w
        = \pi\Bigl(\rho^2 -
        \sum_{n\geq1}n\,\abs{b_n}^2\rho^{-2n}\Bigr) :
        \]
        代入 $\bar g$ 与 $g'$ 在
        $C_\rho$ 上的洛朗级数，并逐项积分（正规
        收敛；只有频率为零的乘积项
        保留下来）。
  \item 令 $\rho \to 1^+$，得到\emph{面积
        定理}：
        \[
        \sum_{n\geq1}n\,\abs{b_n}^2 \;\leq\; 1 .
        \]
        特别地，$\abs{b_1} \leq 1$。何时 $\abs{b_1}
        = 1$ 成立？
\end{enumerate}

\medskip
\noindent\textbf{第 II 部分 --- 比伯巴赫不等式 $\abs{a_2} \leq
2$。} 设 $f = z + a_2z^2 + \cdots \in \mathcal S$。
\begin{enumerate}[resume]
  \item 证明 $f(z^2)/z^2$ 在
        $\mathbb D$ 上全纯且无零点，并存在全纯平方根
        $\varphi$，使 $\varphi(0) = 1$；记 $h(z) =
        z\varphi(z^2)$，从而 $h(z)^2 = f(z^2)$。证明
        $h$ 是\emph{奇}单叶函数，在 $\mathbb
        D$ 上，展开式为 $h(z) = z + \frac{a_2}2z^3 +
        \cdots$。\emph{（单射性：$h(z)^2 = h(z')^2$
        推出 $z^2 = z'^2$；利用奇性完成证明。）}
  \item 对 $g(w) = 1/h(1/w) = w -
        \frac{a_2}{2w} + \cdots$ 在 $\{\abs w > 1\}$ 上应用面积定理
        （验证单叶性与展开式），并得出
        $\abs{a_2} \leq 2$。
  \item 证明\emph{柯贝函数}
        \[
        k(z) = \frac{z}{(1 - z)^2} = \sum_{n\geq1}n\,z^n
        \]
        属于 $\mathcal S$，有 $a_2 = 2$，并将
        $\mathbb D$ 映满 $\C\setminus\intoc{-\infty}{-\frac14}$
        \emph{（写成 $k = \frac14\bigl[\bigl(\frac{1+z}{1 -
        z}\bigr)^2 - 1\bigr]$，并追踪像）}：以下所有
        不等式都是最优的。
\end{enumerate}

\medskip
\noindent\textbf{第 III 部分 --- 四分之一定理。}
\begin{enumerate}[resume]
  \item 设 $f \in \mathcal S$，且 $c \notin f(\mathbb D)$。
        证明：
        \[
        F(z) = \frac{c\,f(z)}{c - f(z)}
        \]
        属于 $\mathcal S$，并计算其第二个系数
        ：$A_2 = a_2 + \frac1c$。
  \item 对 $f$ 与 $F$ 分别应用比伯巴赫不等式，得出
        $\abs{\frac1c} \leq \abs{A_2} + \abs{a_2} \leq 4$，
        即 $\abs c \geq \frac14$。\emph{被略去的每个值的模都
        $\geq \frac14$}：$f(\mathbb D) \supseteq
        D(0, \frac14)$ --- 柯贝四分之一定理。在柯贝函数上验证
        其最优性。
  \item 由此得到黎曼映射的一个定量估计：若
        $\varphi \colon \Omega \to \mathbb D$ 为
        由\cref{thm:b3:conformal:rmt}给出的、在 $z_0$ 处的黎曼映射，则
        \[
        \frac{d\bigl(z_0, \partial\Omega\bigr)}{4}
        \;\leq\; \frac1{\varphi'(z_0)} \;\leq\;
        4\,d\bigl(z_0, \partial\Omega\bigr)
        \]
        --- 至少证明左边的不等式，方法是对
        适当归一化的 $\varphi^{-1}$ 应用柯贝定理，右边
        的不等式则对圆盘上的 $\varphi$ 用施瓦茨引理，圆盘取
        $D(z_0, d) \subseteq \Omega$。
\end{enumerate}

\medskip
\noindent\textbf{第 IV 部分 --- 展望。}
\begin{enumerate}[resume]
  \item 比伯巴赫于 1916 年猜想 $\abs{a_n} \leq n$ 对
        所有 $n$ 成立，且等号仅对柯贝函数的
        旋转取到。德布朗热于 1985 年证明了它。请手算
        验证该猜想对柯贝函数及其
        旋转 $\eu^{-\iu\theta}k(\eu^{\iu\theta}z)$ 成立。然后
        将第 4 问的展开再推进一项，记
        $h(z) = z + \frac{a_2}2z^3 + c_5z^5 + \cdots$，证明
        $c_5 = \frac{a_3}2 - \frac{a_2^2}8$，并
        \[
        g(w) = w - \frac{a_2}{2}\,w^{-1} +
        \Bigl(\frac{3a_2^2}8 - \frac{a_3}2\Bigr)w^{-3} +
        \cdots ,
        \]
        于是由面积定理得到加强的不等式
        $\bigl|\frac{a_2}2\bigr|^2 + 3\bigl|\frac{3a_2^2}8 -
        \frac{a_3}2\bigr|^2 \leq 1$。在柯贝函数上验证
        此不等式（$a_2 = 2$，$a_3 = 3$）。
\end{enumerate}

\medskip
\noindent\textbf{第 V 部分 --- 畸变定理。}
比伯巴赫不等式经
圆盘上的自同构搬运，处处控制 $f'$。取定 $f \in \mathcal
S$。
\begin{enumerate}[resume]
  \item （柯贝变换）设 $z_0 \in \mathbb D$，令
        $\varphi(z) = \frac{z + z_0}{1 + \bar z_0z}$，它是圆盘
        自同构（见\cref{thm:b3:conformal:autdisc}），满足
        $\varphi(0) = z_0$。证明
        \[
        F(z) = \frac{f(\varphi(z)) - f(z_0)}
        {f'(z_0)\,\bigl(1 - \abs{z_0}^2\bigr)}
        \]
        属于 $\mathcal S$ \emph{（单叶性可以继承；
        计算 $\varphi'(0) = 1 - \abs{z_0}^2$，
        并检验归一化条件；回忆 $f' \neq 0$ 对
        单射 $f$，
        见\cref{def:b3:conformal:conformal}）}。
  \item 计算 $A_2 = \frac12F''(0)$，
        即 $F$ 的第二个系数：
        \[
        A_2 = \frac12\Bigl[\bigl(1 - \abs{z_0}^2\bigr)
        \frac{f''(z_0)}{f'(z_0)} - 2\bar z_0\Bigr] ,
        \]
        由比伯巴赫不等式（第 4 题）可推得，对 $z =
        r\eu^{\iu\theta}$，\emph{基本
        不等式}：
        \[
        \Bigl|\,z\,\frac{f''(z)}{f'(z)} -
        \frac{2r^2}{1 - r^2}\Bigr|
        \;\leq\; \frac{4r}{1 - r^2} .
        \]
  \item 取实部：
        \[
        \frac{2r^2 - 4r}{1 - r^2} \;\leq\;
        \operatorname{Re}\Bigl(z\,\frac{f''(z)}{f'(z)}\Bigr)
        \;\leq\; \frac{2r^2 + 4r}{1 - r^2} .
        \]
  \item 证明 $\frac{\dd}{\dd t}\log\bigl|
        f'(t\eu^{\iu\theta})\bigr| = \frac1t
        \operatorname{Re}\bigl(z\frac{f''(z)}{f'(z)}\bigr)$
        在 $z = t\eu^{\iu\theta}$ 处 \emph{（对实变量的非零
        $\mathcal C^1$ 函数 $g$，
        $\frac{\dd}{\dd t}\log\abs g =
        \operatorname{Re}(g'/g)$）}，并沿射线对第
        12 题的界积分，以得到
        \emph{畸变定理}：
        \[
        \frac{1 - r}{(1 + r)^3} \;\leq\; \abs{f'(z)}
        \;\leq\; \frac{1 + r}{(1 - r)^3},
        \qquad \abs z = r .
        \]
  \item 推出\emph{增长定理}：
        \[
        \frac{r}{(1 + r)^2} \;\leq\; \abs{f(z)} \;\leq\;
        \frac{r}{(1 - r)^2},
        \qquad \abs z = r
        \]
        \emph{（上界：对 $f'$ 沿线段
        $[0, z]$ 积分；下界：若 $\abs{f(z)} < \frac14$，
        线段 $[0, f(z)]$ 位于 $f(\mathbb D)$ 中，由
        第 7 题；用 $f^{-1}$ 拉回 ---
        全纯（由
        \cref{cor:b3:residues:openmapping}） --- 并估计
        $\abs{f(z)} = \int_\gamma\abs{f'(\zeta)}
        \,\abs{\dd\zeta} \geq \int_0^r\frac{1 - t}{(1 +
        t)^3}\,\dd t$，利用 $\abs{\dd\zeta} \geq
        \dd\abs\zeta$）}。
  \item 验证柯贝函数在
        全部四个界中都取等号，上界在 $z = r$ 处取等，
        下界在 $z = -r$ 处取等：$k$ 同时是
        扩张性最强、并在对径点处收缩性最强的
        $\mathcal S$ 中的函数。
\end{enumerate}

\medskip
\noindent\textbf{第 VI 部分 --- 极值刚性。}在迄今的每一个不等式中，
等号唯一确定柯贝函数，
        在相差一个旋转的意义下。先证明之，再加以应用。
\begin{enumerate}[resume]
  \item 设 $f \in \mathcal S$ 满足 $\abs{a_2} = 2$。
        沿第 2--4 题追踪等号：面积
        定理给出 $g(w) = w + b_0 + \eu^{\iu\alpha}/w$；
        由 $h$ 为奇函数可知 $g$ 亦为奇函数，故 $b_0 = 0$；求逆
        求出 $h$，再求出 $f$，并由此得出
        \[
        f(z) = \eu^{-\iu\theta}k\bigl(\eu^{\iu\theta}z\bigr)
        \quad\text{ 其中 } \eu^{\iu\theta} =
        -\eu^{\iu\alpha} :
        \]
        柯贝函数的旋转
        是 $\mathcal S$ 中满足 $\abs{a_2} = 2$ 的唯一函数。
  \item 证明：若 $f \in \mathcal S$ 不取值 $c$，
        且 $\abs c = \frac14$ 恰好成立，则 $f$ 为
        柯贝函数的旋转，并确定 $c =
        -\eu^{-\iu\theta}/4$ \emph{（追踪
        等号，经第 7 题的不等式链 $4 = \abs{1/c} = \abs{A_2 - a_2}
        \leq \abs{A_2} + \abs{a_2} \leq 4$）}：四分之一定理的
        常数只在极值
        族中取到。
  \item （用简便方法估计系数）结合增长定理
        与柯西估计
        （见\cref{thm:b3:holomorphic:analytic}）在圆周
        $\abs z = 1 - \frac1n$ 上，证明
        \[
        \abs{a_n} \;\leq\; \eu\,n^2 \qquad (n \geq 2) .
        \]
        （De Branges, 1985：$\abs{a_n} \leq n$；因子
        $\eu n$ 是使用初等工具的代价。）
  \item （子圆盘的覆盖）证明：对每个 $0 < r <
        1$,
        \[
        f\bigl(D(0, r)\bigr) \supseteq
        D\Bigl(0, \frac{r}{(1 + r)^2}\Bigr),
        \]
        对柯贝函数是最优的，并
        令 $r \to 1^-$ 即得四分之一定理。\emph{（开像
        $f(D(0,r))$ 的边界点位于
        $f(\partial D(0,r))$，因而其模 $\geq
        \frac{r}{(1+r)^2}$，由第 14 题；
        从 $0$ 到一个模更小的未覆盖点的线段，
        必定要穿过该边界。）}
  \item （每点处的柯贝估计）设 $f$ 为
        $\mathbb D$ 上的单叶函数（不必归一化），且 $z_0
        \in \mathbb D$。证明
        \[
        \tfrac14\bigl(1 - \abs{z_0}^2\bigr)\abs{f'(z_0)}
        \;\leq\;
        d\bigl(f(z_0), \partial f(\mathbb D)\bigr)
        \;\leq\;
        \bigl(1 - \abs{z_0}^2\bigr)\abs{f'(z_0)}
        \]
        \emph{（左侧：对第 10 题的柯贝
        变换应用四分之一定理；右侧：施瓦茨
        引理（见\cref{thm:b3:conformal:schwarz}）应用于
        $\psi^{-1}\circ\hat g$，其中 $\hat g(w) =
        f^{-1}\bigl(f(z_0) + dw\bigr)$，$d$ 为距离，
        且 $\psi$ 为将 $0$ 映到
        $z_0$ 的圆盘自同构）}。为何 $\partial f(\mathbb D)$ 非空？
  \item 检验第 20 题在 $f = k$，$z_0 = r \in
        \intoo01$ 时的结论：计算 $d\bigl(k(r), \partial k(\mathbb
        D)\bigr) = \frac{(1+r)^2}{4(1-r)^2}$，并验证
        左侧不等式实际上是\emph{等式}：柯贝
        函数使自身定理在
        $\intoo01$ 的每一点上都取等号。
  \item （总结）用十行以内说明：哪一个单一原理
        构成面积定理的基础，以及比伯巴赫不等式、
        四分之一定理、畸变、增长与覆盖如何
        皆由其导出？可与
        施瓦茨--皮克体系相比较（见
        \cref{thm:b3:conformal:schwarz}）：二者都由一个内部
        不等式使整个几何刚性化，
        且极值函数在相差一个旋转的意义下唯一。

\end{enumerate}
\medskip
\noindent\textbf{第 VII 部分 --- 紧性、逆映射与
现实检验。}
\begin{enumerate}[resume]
  \item 证明类 $\mathcal S$ 关于局部一致收敛是\emph{紧}的：
        即由增长定理，它局部有界，
        从而正规
        （\cref{thm:b3:conformal:montel}）；并且
        $\mathcal S$ 中元素的局部一致极限
        仍属于 $\mathcal S$ \emph{（归一化条件经
        导数的魏尔斯特拉斯收敛传到极限；
        单射性由赫尔维茨定理保持（
        见\cref{exo:b3:residues:8}，极限非常值）}。
        这对于比伯巴赫这类
        极值问题有何意义？
  \item 对 $f \in \mathcal S$，令 $g = f^{-1}$，
        在 $0$ 附近有定义。证明 $g(w) = w - a_2w^2 + O(w^3)$，
        于是逆函数的第二系数满足同样的最佳
        估计 $\abs{A_2} = \abs{a_2} \leq 2$，等号恰在
        旋转后的柯贝函数时成立。
  \item 确定对哪些 $a \in \C$，多项式 $f(z)
        = z + az^2$ 属于 $\mathcal S$：证明 $f$ 在
        $\mathbb D$ 上单射当且仅当 $\abs a \leq \frac12$
        \emph{（因式分解 $f(z_1) - f(z_2)$）}。由此得出：
        对二次多项式，真正的系数界为
        $\frac12$，比伯巴赫的界 $2$ 小四倍
        --- $\mathcal S$ 的极值函数确实是
        真正的超越对象，没有任何多项式
        能与之接近。
\end{enumerate}
\end{problem}
